%!TEX program = xelatex
\documentclass[11pt]{ctexart}

\usepackage[margin=1in]{geometry}
\usepackage{amsmath, amssymb, amsthm, bm}
\usepackage{hyperref}
\usepackage{xcolor}
\usepackage{enumitem}
\usepackage{iftex}

% Robust CJK font fallback across systems
\ifXeTeX
  \newif\ifcjkfontset
  \cjkfontsetfalse
  \IfFontExistsTF{Noto Serif CJK SC}{\setCJKmainfont[AutoFakeBold=2,AutoFakeSlant=true]{Noto Serif CJK SC}\cjkfontsettrue}{}
  \ifcjkfontset\else
    \IfFontExistsTF{Source Han Serif SC}{\setCJKmainfont[AutoFakeBold=2,AutoFakeSlant=true]{Source Han Serif SC}\cjkfontsettrue}{}
  \fi
  \ifcjkfontset\else
    \IfFontExistsTF{Songti SC}{\setCJKmainfont[AutoFakeBold=2,AutoFakeSlant=true]{Songti SC}\cjkfontsettrue}{}
  \fi
  \ifcjkfontset\else
    \IfFontExistsTF{SimSun}{\setCJKmainfont[AutoFakeBold=2,AutoFakeSlant=true]{SimSun}\cjkfontsettrue}{}
  \fi
  \ifcjkfontset\else
    \usepackage{iftex}

% Robust CJK font fallback across systems
\ifXeTeX
  \newif\ifcjkfontset
  \cjkfontsetfalse
  \IfFontExistsTF{Noto Serif CJK SC}{\setCJKmainfont[AutoFakeBold=2,AutoFakeSlant=true]{Noto Serif CJK SC}\cjkfontsettrue}{}
  \ifcjkfontset\else
    \IfFontExistsTF{Source Han Serif SC}{\setCJKmainfont[AutoFakeBold=2,AutoFakeSlant=true]{Source Han Serif SC}\cjkfontsettrue}{}
  \fi
  \ifcjkfontset\else
    \IfFontExistsTF{Songti SC}{\setCJKmainfont[AutoFakeBold=2,AutoFakeSlant=true]{Songti SC}\cjkfontsettrue}{}
  \fi
  \ifcjkfontset\else
    \IfFontExistsTF{SimSun}{\setCJKmainfont[AutoFakeBold=2,AutoFakeSlant=true]{SimSun}\cjkfontsettrue}{}
  \fi
  \ifcjkfontset\else
    \setCJKmainfont[AutoFakeBold=2,AutoFakeSlant=true]{FandolSong-Regular}
  \fi
\else
  \PackageError{plan_discussion_further}{Please compile with XeLaTeX}{Use xelatex instead of pdflatex/lualatex.}
\fi  \fi
\else
  \PackageError{plan_discussion_further}{Please compile with XeLaTeX}{Use xelatex instead of pdflatex/lualatex.}
\fi

\newtheorem{theorem}{定理}
\newtheorem{proposition}{命题}
\newtheorem{corollary}{推论}
\newtheorem{definition}{定义}
\newtheorem{remark}{注记}
\newtheorem{assumption}{假设}

\setlength{\parindent}{0pt}
\setlength{\parskip}{0.8\baselineskip}

\title{未知后续损失下的 Rayleigh 下界推导\\
\large 当 $z = Wx$ 后面不再假设 softmax + cross-entropy 时，理论需要什么条件}
\author{Codex}
\date{\today}

\begin{document}
\maketitle

\section*{0. 写作目标}

本文讨论如下更一般的问题：

\[
z = Wx, \qquad W \in \mathbb{R}^{d_{out} \times d_{in}}, \quad x \in \mathbb{R}^{d_{in}}, \quad z \in \mathbb{R}^{d_{out}}.
\]

在 $z$ 之后，我们\textbf{不再假设}接 softmax，也\textbf{不再假设}损失是 cross-entropy。我们只知道对每个样本存在某个标量损失
\[
\ell(z; \xi),
\]
其中 $\xi$ 表示一切后续监督或环境信号；其具体形式未知。

我们要回答三个问题：

\begin{enumerate}[label=\arabic*.]
    \item 在这种一般情形下，原有的 Rayleigh 商推导是否还能成立？
    \item 若能成立，究竟需要对 $H_z := \nabla_z^2 \ell(z;\xi)$ 施加哪些条件？
    \item 哪些条件只是为了\textbf{代数推导成立}，哪些条件是为了把结果解释成\textbf{曲率 / sharpness 屏障}？
\end{enumerate}

结论先说在前面：

\begin{itemize}
    \item \textbf{代数推导本身并不依赖 softmax 或 cross-entropy。}
    \item 真正核心的是：$z=Wx$ 必须是线性的，且 $H_z$ 必须作为一个对称二阶算子存在并满足基本可积性。
    \item 如果还想把结论解释成“正曲率”“sharpness”“Gauss-Newton 屏障”，则必须进一步假设 $H_z \succeq 0$（至少在我们关心的方向上非负，最好几乎处处半正定）。
    \item 如果你\textbf{只有一阶反传信号}而完全不知道损失形式，那么 $H_z$ 本身并不能从单次反传中被唯一恢复；此时理论只能写成一个\textbf{条件性结论}，或者转而研究某个\textbf{代理矩阵}。
\end{itemize}

\section*{1. 一般设定与符号}

令样本 $(x,\xi)$ 服从某个分布 $\mathcal{D}$。如果你研究的是某个 expert 的条件子分布，也可以把 $\mathcal{D}$ 替换成 $\mathcal{D}_e$；下述推导不依赖 MoE 结构本身。

\begin{assumption}[未知损失的最小解析假设]\label{ass:basic_loss}
对几乎处处的样本 $(x,\xi)$，函数 $\ell(\cdot;\xi)$ 在 $z=Wx$ 的邻域内二阶可微，并定义
\[
H_z := \nabla_z^2 \ell(z;\xi) \in \mathbb{R}^{d_{out}\times d_{out}}.
\]
此外假设
\[
\mathbb{E}\big[\|x\|_2^2 \, \|H_z\|_{\mathrm{op}}\big] < \infty.
\]
\end{assumption}

\begin{remark}
由于 $H_z$ 是标量函数对 $z$ 的 Hessian，只要二阶可微，就自动满足
\[
H_z = H_z^\top.
\]
因此“对称性”不是额外假设，而是 Hessian 作为 Hessian 的内在性质。
\end{remark}

\begin{definition}[输出空间平均 Hessian 与有效曲率]
定义
\[
\Sigma_H := \mathbb{E}[H_z].
\]
由于 $\Sigma_H$ 是有限维实对称矩阵，最大特征值存在。记
\[
\gamma_{\mathrm{eff}} := \lambda_{\max}(\Sigma_H),
\]
并取任意对应单位特征向量 $u_H \in \mathbb{R}^{d_{out}}$，满足
\[
\|u_H\|_2 = 1, \qquad \Sigma_H u_H = \gamma_{\mathrm{eff}} u_H.
\]
对单个样本，定义方向二阶量
\[
h := u_H^\top H_z u_H.
\]
\end{definition}

\begin{remark}
到这里为止，我们\textbf{完全没有使用 softmax，也没有使用 cross-entropy}。唯一对象是未知损失的输出空间 Hessian $H_z$。
\end{remark}

\section*{2. 线性块的精确二阶结构}

下面是整个框架最核心、也最不依赖具体损失形式的一步。

\begin{proposition}[当 $z=Wx$ 时，参数块 Hessian 具有 Kronecker 结构]\label{prop:exact_hessian}
令 $w := \mathrm{vec}(W) \in \mathbb{R}^{d_{out}d_{in}}$。在假设 \ref{ass:basic_loss} 下，对单个样本有
\[
\nabla_w^2 \ell(z;\xi) = xx^\top \otimes H_z.
\]
因此其期望参数块 Hessian 为
\[
\mathcal{H} := \mathbb{E}\big[xx^\top \otimes H_z\big].
\]
\end{proposition}

\begin{proof}
将 $z = Wx$ 写成关于 $w = \mathrm{vec}(W)$ 的线性映射。记
\[
J := \frac{\partial z}{\partial w} = x^\top \otimes I_{d_{out}}.
\]
因为 $z$ 对 $w$ 是\textbf{线性的}，所以 $z$ 关于 $w$ 的二阶导数恒为零。对标量函数 $\ell(z;\xi)$ 应用链式法则，
\[
\nabla_w \ell = J^\top \nabla_z \ell,
\]
再次求导得
\[
\nabla_w^2 \ell = J^\top H_z J.
\]
将 $J = x^\top \otimes I_{d_{out}}$ 代入，并利用 Kronecker 乘积恒等式，
\[
J^\top H_z J
= (x \otimes I_{d_{out}}) H_z (x^\top \otimes I_{d_{out}})
= xx^\top \otimes H_z.
\]
最后对样本取期望即可得到
\[
\mathcal{H} = \mathbb{E}[xx^\top \otimes H_z].
\]
\end{proof}

\begin{remark}[这一块是“精确 Hessian”还是“Gauss-Newton”？]
在本文的一般设定里，$\mathcal{H}$ 是\textbf{精确参数块 Hessian}，因为我们直接对未知标量损失 $\ell(z;\xi)$ 求二阶导。

只有在 $H_z \succeq 0$ 时，人们才常把
\[
\mathcal{H} = \mathbb{E}[xx^\top \otimes H_z]
\]
理解为 generalized Gauss-Newton / Fisher 型正曲率块。若 $H_z$ 可正可负，则此处最好称其为\textbf{输出空间 Hessian 诱导的精确块 Hessian}，而不是默认称作 GN 块。
\end{remark}

\section*{3. 激活分解与一维投影}

下面沿用原理论的统计分解，但不再假设任何特定损失形式。

\begin{assumption}[激活均值分解]\label{ass:mean_decomp}
假设
\[
x = \mu + \tilde{x}, \qquad \mu := \mathbb{E}[x], \qquad \mathbb{E}[\tilde{x}] = 0.
\]
并且 $\mathbb{E}\|x\|_2^2 < \infty$。
\end{assumption}

\begin{assumption}[非退化均值方向]\label{ass:mu_nonzero}
暂时假设 $\mu \neq 0$，并定义
\[
u_\mu := \frac{\mu}{\|\mu\|_2}.
\]
记沿均值方向的一维残差为
\[
r := u_\mu^\top \tilde{x}.
\]
定义
\[
\sigma_\mu^2 := \mathrm{Var}(r).
\]
\end{assumption}

\begin{remark}[当 $\mu = 0$ 时怎么办]
若 $\mu = 0$，则“均值方向”本身不存在。此时原理论中的 $u_\mu$ 不再有规范定义。

一种严格的做法是：
\begin{itemize}
    \item 要么把定理写成“沿任意给定单位向量 $a$ 的版本”，把 $u_\mu$ 全部替换为 $a$；
    \item 要么额外引入一个数据驱动锚定方向（例如激活协方差的主方向），再沿该方向写同样的推导。
\end{itemize}

因此，下文的“均值方向版本”严格地说建立在 $\mu \neq 0$ 之上。
\end{remark}

\section*{4. 一般情形下的 Rayleigh 商恒等式}

\begin{theorem}[未知损失下的通用 Rayleigh 恒等式]\label{thm:generic_exact_identity}
在假设 \ref{ass:basic_loss}、\ref{ass:mean_decomp}、\ref{ass:mu_nonzero} 下，定义
\[
v := u_\mu \otimes u_H.
\]
则有
\[
\lambda_{\max}(\mathcal{H}) \ge v^\top \mathcal{H} v,
\]
且
\begin{equation}\label{eq:generic_exact_identity}
v^\top \mathcal{H} v
= \gamma_{\mathrm{eff}}\big(\|\mu\|_2^2 + \sigma_\mu^2\big)
 2\|\mu\|_2 \,\mathrm{Cov}(r,h)
 \mathrm{Cov}(r^2,h).
\end{equation}
从而
\begin{equation}\label{eq:generic_exact_rayleigh_lower}
\lambda_{\max}(\mathcal{H})
\ge
\gamma_{\mathrm{eff}}\big(\|\mu\|_2^2 + \sigma_\mu^2\big)
 2\|\mu\|_2 \,\mathrm{Cov}(r,h)
 \mathrm{Cov}(r^2,h).
\end{equation}
\end{theorem}

\begin{proof}
由 Rayleigh 商原理，
\[
\lambda_{\max}(\mathcal{H}) \ge v^\top \mathcal{H} v.
\]
由于 $\|u_\mu\|_2 = \|u_H\|_2 = 1$，所以 $\|v\|_2=1$。

利用命题 \ref{prop:exact_hessian}，
\[
\mathcal{H} = \mathbb{E}[xx^\top \otimes H_z].
\]
于是
\begin{align*}
v^\top \mathcal{H} v
&= (u_\mu \otimes u_H)^\top \, \mathbb{E}[xx^\top \otimes H_z] \, (u_\mu \otimes u_H) \\
&= \mathbb{E}\big[(u_\mu^\top x)^2 (u_H^\top H_z u_H)\big] \\
&= \mathbb{E}\big[(u_\mu^\top x)^2 h\big].
\end{align*}

由 $x=\mu+\tilde{x}$ 及 $u_\mu^\top \mu = \|\mu\|_2$，
\[
u_\mu^\top x = \|\mu\|_2 + r.
\]
因此
\begin{align*}
v^\top \mathcal{H} v
&= \mathbb{E}\big[(\|\mu\|_2 + r)^2 h\big] \\
&= \|\mu\|_2^2 \mathbb{E}[h]
 + 2\|\mu\|_2 \mathbb{E}[rh]
 + \mathbb{E}[r^2 h].
\end{align*}

下面分别展开三项。

\paragraph{第一项}
由 $u_H$ 是 $\Sigma_H = \mathbb{E}[H_z]$ 的最大特征向量，
\[
\mathbb{E}[h]
= u_H^\top \mathbb{E}[H_z] u_H
= u_H^\top \Sigma_H u_H
= \gamma_{\mathrm{eff}}.
\]
故第一项为
\[
\gamma_{\mathrm{eff}} \|\mu\|_2^2.
\]

\paragraph{第二项}
由于 $\mathbb{E}[r]=0$，
\[
\mathbb{E}[rh] = \mathrm{Cov}(r,h).
\]
故第二项为
\[
2\|\mu\|_2 \,\mathrm{Cov}(r,h).
\]

\paragraph{第三项}
由协方差恒等式，
\[
\mathbb{E}[r^2 h]
= \mathbb{E}[r^2]\mathbb{E}[h] + \mathrm{Cov}(r^2,h)
= \sigma_\mu^2 \gamma_{\mathrm{eff}} + \mathrm{Cov}(r^2,h).
\]

将三项合并，即得
\[
v^\top \mathcal{H} v
= \gamma_{\mathrm{eff}}\big(\|\mu\|_2^2+\sigma_\mu^2\big)
 2\|\mu\|_2\,\mathrm{Cov}(r,h)
 \mathrm{Cov}(r^2,h).
\]
这就证明了 \eqref{eq:generic_exact_identity} 与 \eqref{eq:generic_exact_rayleigh_lower}。
\end{proof}

\begin{remark}[这是最一般、也最本质的形式]
定理 \ref{thm:generic_exact_identity} 是整个框架在“未知损失”情形下的\textbf{核心形式}。到这一步为止，既没有用 softmax，也没有用 cross-entropy，更没有用 $H_z \succeq 0$。

真正被用到的只有：
\begin{enumerate}[label=(\alph*)]
    \item $z = Wx$ 的线性性；
    \item $H_z$ 存在且可积；
    \item $u_H$ 是 $\mathbb{E}[H_z]$ 的最大特征向量；
    \item 一维投影 $r = u_\mu^\top \tilde{x}$ 的统计量存在。
\end{enumerate}
\end{remark}

\section*{5. 如何得到与原文同型的下界}

原文的界是
\[
\lambda_{\max} \ge \gamma_{\mathrm{eff}}(\|\mu\|^2 + \sigma_\mu^2)
- 2\|\mu\||\rho|\sigma_\mu \sigma_h
- \epsilon_{\mathrm{var}}.
\]

这一形式并不依赖 CE；它只是把定理 \ref{thm:generic_exact_identity} 中的两个协方差项做了保守下界。

\begin{assumption}[二阶统计存在]\label{ass:moment_finite}
假设
\[
\mathbb{E}[h^2] < \infty, \qquad \mathbb{E}[r^4] < \infty.
\]
于是 $\sigma_h^2 := \mathrm{Var}(h)$ 存在，且 $\mathrm{Cov}(r,h)$、$\mathrm{Cov}(r^2,h)$ 都有定义。
\end{assumption}

\begin{definition}[一阶相关与二阶惩罚]
在假设 \ref{ass:moment_finite} 下，定义
\[
\rho := \mathrm{Corr}(r,h),
\qquad
\sigma_h^2 := \mathrm{Var}(h),
\]
以及
\[
\epsilon_{\mathrm{var}} := \max\bigl(0,\,-\mathrm{Cov}(r^2,h)\bigr).
\]
于是自动有
\[
\mathrm{Cov}(r^2,h) \ge -\epsilon_{\mathrm{var}}.
\]
\end{definition}

\begin{corollary}[未知损失下的同型保守下界]\label{cor:generic_bound}
在定理 \ref{thm:generic_exact_identity} 及假设 \ref{ass:moment_finite} 下，
\begin{equation}\label{eq:generic_bound}
\lambda_{\max}(\mathcal{H})
\ge
\gamma_{\mathrm{eff}}\big(\|\mu\|_2^2 + \sigma_\mu^2\big)
- 2\|\mu\|_2 |\rho| \sigma_\mu \sigma_h
- \epsilon_{\mathrm{var}}.
\end{equation}
\end{corollary}

\begin{proof}
由定理 \ref{thm:generic_exact_identity}，
\[
\lambda_{\max}(\mathcal{H})
\ge
\gamma_{\mathrm{eff}}\big(\|\mu\|_2^2 + \sigma_\mu^2\big)
 2\|\mu\|_2 \mathrm{Cov}(r,h)
 \mathrm{Cov}(r^2,h).
\]

又由相关系数定义，
\[
\mathrm{Cov}(r,h) = \rho \sigma_\mu \sigma_h,
\]
因此
\[
2\|\mu\|_2 \mathrm{Cov}(r,h)
\ge
-2\|\mu\|_2 |\rho| \sigma_\mu \sigma_h.
\]

同时按 $\epsilon_{\mathrm{var}}$ 的定义，
\[
\mathrm{Cov}(r^2,h) \ge -\epsilon_{\mathrm{var}}.
\]
代回即可得到 \eqref{eq:generic_bound}。
\end{proof}

\begin{remark}
所以，\textbf{原公式的形式并不依赖 cross-entropy}。它只依赖于：
\begin{itemize}
    \item 你能定义一个输出空间 Hessian $H_z$；
    \item 这个 Hessian 的平均最大特征方向 $u_H$ 存在；
    \item 沿 $u_H$ 的样本方向量 $h = u_H^\top H_z u_H$ 统计上可控。
\end{itemize}

cross-entropy 只是一个让这些量\textbf{自动具备良好性质}的特殊例子，而不是这套代数结构的必要前提。
\end{remark}

\section*{6. 到底必须对 $H_z$ 加什么条件？}

这是本文最关键的问题。下面严格区分两个层次。

\subsection*{6.1 只为了让推导成立：最小条件}

如果你的目标只是让定理 \ref{thm:generic_exact_identity} 和推论 \ref{cor:generic_bound} \textbf{作为代数结论成立}，则对 $H_z$ 需要的最小条件是：

\begin{enumerate}[label=\textbf{(M\arabic*)}]
    \item \textbf{Hessian 存在：} $\ell(\cdot;\xi)$ 在关注区域对 $z$ 二阶可微，因此 $H_z = \nabla_z^2 \ell(z;\xi)$ 存在。
    \item \textbf{对称性：} 由于 $H_z$ 是 Hessian，故 $H_z = H_z^\top$。
    \item \textbf{可积性：} $\mathbb{E}[\|x\|_2^2 \|H_z\|_{\mathrm{op}}] < \infty$，从而
    \[
    \mathcal{H} = \mathbb{E}[xx^\top \otimes H_z]
    \]
    有意义。
    \item \textbf{方向二阶矩存在：} 对 $u_H$ 定义的
    \[
    h = u_H^\top H_z u_H
    \]
    需要满足 $\mathbb{E}[h^2] < \infty$，以便定义 $\sigma_h$ 与 $\rho$。
    \item \textbf{四阶激活矩存在：} 至少要有 $\mathbb{E}[r^4] < \infty$，以便 $\mathrm{Cov}(r^2,h)$ 存在。
\end{enumerate}

在这些最小条件下，\textbf{推导已经成立}，即：
\begin{itemize}
    \item 精确参数块 Hessian 具有 Kronecker 结构；
    \item Rayleigh 商恒等式成立；
    \item 同型的保守下界成立。
\end{itemize}

\textbf{重要：这里完全不需要 $H_z \succeq 0$。}

\subsection*{6.2 若还想把它解释成“正曲率 / sharpness 屏障”}

如果你不只想要一个代数不等式，而是想说：

\begin{quote}
“这个下界描述了训练中的正曲率、sharpness、稳定性边界，因而构成某种 curvature barrier。”
\end{quote}

那么对 $H_z$ 需要加强条件：

\begin{enumerate}[label=\textbf{(C\arabic*)}]
    \item \textbf{输出空间半正定：}
    \[
    H_z \succeq 0 \qquad \text{几乎处处成立。}
    \]
    这是最强、也最干净的条件。

    \item \textbf{非退化有效曲率：}
    \[
    \gamma_{\mathrm{eff}} = \lambda_{\max}(\Sigma_H) > 0.
    \]
    否则主导项本身可能为零或负。

    \item \textbf{修正项不能压倒主项：}
    至少在你关心的训练阶段，需要
    \[
    \gamma_{\mathrm{eff}}(\|\mu\|_2^2 + \sigma_\mu^2)
    >
    2\|\mu\|_2 |\rho| \sigma_\mu \sigma_h + \epsilon_{\mathrm{var}}.
    \]
    只有这样，下界才是严格正的，从而真能被解释为“正曲率屏障”。
\end{enumerate}

\begin{remark}[为什么 $H_z \succeq 0$ 很关键]
如果 $H_z$ 几乎处处半正定，则：
\[
h = u_H^\top H_z u_H \ge 0,
\qquad
xx^\top \otimes H_z \succeq 0,
\qquad
\mathcal{H} \succeq 0.
\]
于是 $\lambda_{\max}(\mathcal{H})$ 的确是一个\textbf{非负曲率强度}，而不是“一个可能带有非凸性混入的最大特征值”。

反之，若 $H_z$ 是不定的，则下界 \eqref{eq:generic_bound} 仍然代数上正确，但此时它应被理解为
\begin{quote}
“未知二阶算子诱导的某个最大特征值下界”
\end{quote}
而不应直接等同于正曲率 sharpness。
\end{remark}

\subsection*{6.3 只有一阶损失信号时，为什么还不够}

如果你知道的仅仅是一次反传给出的
\[
g_z := \nabla_z \ell(z;\xi),
\]
但\textbf{不知道} $\ell$ 的具体形式，也\textbf{没有}对 $z$ 的 Hessian-vector product 接口，那么单凭 $g_z$ 无法唯一恢复 $H_z$。

原因很简单：很多不同的标量函数在同一点可以有\textbf{相同的一阶导}，但\textbf{不同的二阶导}。因此：

\begin{itemize}
    \item 只知道 $g_z$，不足以确定 $H_z$；
    \item 不知道 $H_z$，就无法把理论实例化为一个可计算的数值对象；
    \item 这时理论只能写成：
    \begin{quote}
    “若存在某个满足 (M1)--(M5) 或 (C1)--(C3) 的 $H_z$，则以下结论成立。”
    \end{quote}
\end{itemize}

\begin{remark}[工程上意味着什么]
如果你既没有显式损失形式，也没有对 $z$ 的二阶 oracle，那么你只能走两条路：

\begin{enumerate}[label=\arabic*.]
    \item \textbf{条件理论路线：} 把 $H_z$ 当作抽象二阶算子，只陈述“若其满足某些条件，则结论成立”；
    \item \textbf{代理矩阵路线：} 人为选一个代理 $G_z$ 来代替真实 $H_z$，例如某种 Fisher / softmax Jacobian / 局部 PSD 近似。
\end{enumerate}

但一旦你走代理路线，结论就变成了
\[
\widetilde{\mathcal{H}} := \mathbb{E}[xx^\top \otimes G_z]
\]
的理论，而\textbf{不再是}真实损失 Hessian 的理论。
\end{remark}

\section*{7. softmax + cross-entropy 为什么只是一个特殊且方便的例子}

现在回头看 cross-entropy 情形。若
\[
p = \mathrm{softmax}(z),
\qquad
\ell(z;y) = -\log p_y,
\]
则有著名恒等式
\[
H_z = \nabla_z^2 \ell(z;y) = \mathrm{diag}(p) - pp^\top.
\]

这时有几个额外好处：

\begin{enumerate}[label=\arabic*.]
    \item $H_z$ 有显式闭式表达；
    \item $H_z \succeq 0$ 自动成立；
    \item $h = u_H^\top H_z u_H \ge 0$ 自动成立；
    \item $\mathcal{H} = \mathbb{E}[xx^\top \otimes H_z] \succeq 0$ 自动成立；
    \item “曲率 / sharpness / barrier”的解释自然成立。
\end{enumerate}

所以，softmax + CE 的真正作用不是“让推导第一次成立”，而是：

\begin{quote}
它给了我们一个显式、半正定、易解释、易计算的 $H_z$。
\end{quote}

这正是它在原始文档中显得特别自然的原因。

\section*{8. 最终结论：严格回答“需要对 $H_z$ 加什么条件”}

如果问题是：

\begin{quote}
“当 $z=Wx$ 后面不接 softmax，甚至不知道损失形式时，我们需要人为地对 $H_z$ 加上什么条件，才能成立原来的推导？”
\end{quote}

那么严格答案是：

\subsection*{8.1 为了让\textbf{推导}成立，最小条件是}

\begin{enumerate}[label=\textbf{(A\arabic*)}]
    \item $H_z = \nabla_z^2 \ell(z;\xi)$ 存在；
    \item $H_z$ 是对称矩阵；
    \item $\mathbb{E}[\|x\|_2^2 \|H_z\|_{\mathrm{op}}] < \infty$；
    \item 沿主方向的二阶量 $h = u_H^\top H_z u_H$ 满足 $\mathbb{E}[h^2] < \infty$；
    \item 激活投影残差满足 $\mathbb{E}[r^4] < \infty$；
    \item 若使用“均值方向版本”，还需 $\mu \neq 0$。
\end{enumerate}

这就足以得到本文的通用恒等式 \eqref{eq:generic_exact_identity} 和保守下界 \eqref{eq:generic_bound}。

\subsection*{8.2 为了让它具有\textbf{曲率 / sharpness}解释，还要再加}

\begin{enumerate}[label=\textbf{(B\arabic*)}]
    \item $H_z \succeq 0$ 几乎处处成立；
    \item $\gamma_{\mathrm{eff}} > 0$；
    \item 主项在关心阶段压过耦合修正项。
\end{enumerate}

\subsection*{8.3 如果你连真实 $H_z$ 都不知道}

那就不能宣称你在研究“真实损失曲率”；你最多只能：

\begin{enumerate}[label=\arabic*.]
    \item 写一个关于抽象 $H_z$ 的条件性定理；
    \item 或者换成某个代理矩阵 $G_z$，并明确声明你研究的是\textbf{代理曲率}。
\end{enumerate}

\begin{remark}[一句话总结]
\textbf{原推导真正依赖的不是 softmax，而是“输出空间二阶算子 $H_z$ 的存在、对称性、可积性”。}

\textbf{原推导真正依赖 cross-entropy 的地方，是把这个抽象的 $H_z$ 具体化成了一个天然半正定、显式可算的矩阵 $\mathrm{diag}(p)-pp^\top$。}
\end{remark}

\end{document}
